% !TEX root = ../main.tex
\chapter{仕様・実装・テスト・証明の接続}
\label{chap:spec-impl-test-proof}
\BookIndex{てすと}{テスト}
\BookIndex{しょうめい}{証明}
\BookIndex{ひんしつほしょう}{品質保証}

\begin{chapterabstract}
本章では、仕様、実装、テスト、証明を別々の活動としてではなく、一つの品質保証の流れとして接続します。
Lean 4 による証明は、すべての品質問題を解決する万能薬ではありません。
外部 API、UI、性能、運用環境、障害時挙動などには、通常のテスト、監視、レビューのほうが向いています。
一方、round-trip 性質、不変条件保存、件数保存、処理の意味保存のように、すべての入力について確認したい小さな核は Lean で扱いやすい対象です。
本章の目的は、何を Lean で証明し、何を通常のテストに残し、両者を設計文書の中でどのように結びつけるかを判断できるようにすることです。
\end{chapterabstract}

\begin{learninggoals}
\item Lean で証明する価値が高い性質を見分けられます。
\item 通常のテスト、型、仕様、証明の役割分担を説明できます。
\item 自然言語仕様から、証明可能な小さな核を抽出できます。
\item Lean theorem を設計文書やレビュー項目に対応づける書き方を理解できます。
\end{learninggoals}

\section{仕様・実装・テスト・証明の役割}

ここまでの章では、リファクタリングを等式として扱い、マイグレーションを同型や片方向互換として分類し、API adapter の一貫性を自然性として表しました。
これらはすべて、ソフトウェア上の変更を「何が保存されるべきか」という問いに置き換える作業でした。

しかし実務では、次の点を判断する必要があります。

\begin{itemize}
\item どの性質を Lean で証明するか。
\item どの挙動は単体テストで確認すれば十分か。
\item 外部 API の応答や DB の実挙動を Lean のモデルに含めるか。
\item 証明した theorem を、設計文書やレビューでどのように読むか。
\end{itemize}

この判断を誤ると、Lean の導入は重くなります。
何でも形式化しようとすると、仕様モデルが本番コードより複雑になり、変更するたびに保守できなくなります。
逆に、テストだけに寄せすぎると、すべての入力に対して成り立つべき性質を代表例の確認に落としてしまいます。

Lean は、テストを置き換えるものではありません。
Lean は、テストでは代表例しか見られない性質のうち、小さく、純粋で、壊れると困る核を検査する道具です。
テストは実環境に近い振る舞いを見ます。
証明はモデル化した範囲の全ケースを見ます。
テストと証明は競合するのではなく、役割が違います。

\FigurePlaceholder{fig-ch35-connection-map.png}{0.92}{仕様・実装・テスト・証明}{fig:ch35-connection-map}

図\ref{fig:ch35-connection-map}では、仕様から実装へ直接進む線だけでなく、テストと証明を横に置いています。
テストは「実装が期待例で動くか」を見ます。
証明は「モデル化した仕様がすべての場合に成り立つか」を見ます。
設計文書は、その対応関係と前提条件を人間が読める形で残す場所です。

\section{何を Lean で証明するべきか}

Lean で証明する価値が高いのは、実装のすべてではありません。
価値が高いのは、次の条件を多く満たす部分です。

\begin{itemize}
\item 入力と出力を小さな型として切り出せます。
\item 外部状態に依存しない純粋な関数として表せます。
\item 代表例ではなく、すべての入力について保証したい性質です。
\item 仕様変更やリファクタリングで壊れると被害が大きい性質です。
\item 証明したい性質を、等式、含意、不変条件、round-trip などとして短く書けます。
\end{itemize}

本章では、\textbf{証明可能な核}という言葉を使います。
証明可能な核とは、実務仕様全体のうち、Lean 上で小さな型、関数、述語、定理として表せる中心部分です。
たとえば、料金計算の丸め規則すべてではなく、「この二つの実装は同じモデル関数に一致する」という等式だけを切り出す場合、その等式が証明可能な核になります。

典型的には、次のような性質が Lean に向いています。

\begin{longtable}{p{0.30\linewidth}p{0.61\linewidth}}
\toprule
\textbf{性質} & \textbf{Lean に向く理由} \\
\midrule
round-trip & 変換して戻すと元に戻るという等式で書けます。
スキーマ移行や encoder/decoder に向いています。 \\
不変条件保存 & 操作前に条件を満たすなら、操作後も条件を満たすという含意で書けます。
状態更新の核に向いています。 \\
件数保存 & \texttt{List.map} などの構造保存として書けます。
バッチ移行の基本性質になります。 \\
意味保存 & 変更前後の関数がすべての入力で同じ結果を返すという等式で書けます。
リファクタリングに向いています。 \\
可換性 & adapter と業務ロジックの順序を入れ替えても同じという等式で書けます。
自然性の小さな実務版です。 \\
\bottomrule
\end{longtable}

\subsection{小さな料金計算モデル}

まず、仕様と実装を分ける最小例を見ます。
ここでは、価格、割引、送料を自然数で扱います。
現実の通貨、税率、丸め、クーポン期限、外部決済 API はまだ扱いません。
これらを最初から入れると、証明の核が見えにくくなるためです。

次の Lean コードでは、\texttt{chargeSpec} を仕様モデル、\texttt{chargeImpl} を実装モデルとして分けています。
定理 \texttt{charge\_impl\_matches\_spec} は、実装モデルが仕様モデルと同じ結果を返すことを、任意の入力 \texttt{i} について主張しています。

\begin{listing}[htbp]
\begin{minted}{lean4}
namespace Chapter35

structure PriceInput where
  base : Nat
  discount : Nat
  shipping : Nat

def chargeSpec (i : PriceInput) : Nat :=
  (i.base - i.discount) + i.shipping

def chargeImpl (i : PriceInput) : Nat :=
  let subtotal := i.base - i.discount
  subtotal + i.shipping

theorem charge_impl_matches_spec (i : PriceInput) :
    chargeImpl i = chargeSpec i := by
  rfl
\end{minted}
\caption{\texttt{charge\_impl\_matches\_spec}}
\label{lst:ch35-charge-spec-impl}
\end{listing}

ここで証明しているのは、料金システム全体の正しさではありません。
証明しているのは、\texttt{chargeImpl} という小さな実装モデルが、\texttt{chargeSpec} という小さな仕様モデルと一致することだけです。
しかも、この例の二定義は展開すると同じ式なので、定理だけから \texttt{chargeSpec} が業務仕様として妥当だとは分かりません。
仕様モデルを独立にレビューし、本番実装との対応を保つ作業が別に必要です。

実務では、\texttt{chargeSpec} が仕様書中の数式や業務ルールに対応し、\texttt{chargeImpl} が本番コードから切り出した純粋な核に対応します。
Lean の theorem は、「この核については、実装が仕様モデルからずれていない」というレビュー可能な証拠になります。

\section{何は通常のテストでよいか}

Lean は強力ですが、現実のシステムのすべてを直接扱うには向いていません。
通常のテストに残すべきものも多くあります。

\begin{longtable}{p{0.32\linewidth}p{0.59\linewidth}}
\toprule
\textbf{通常テストが向くもの} & \textbf{理由} \\
\midrule
外部 API との通信 & 相手の実装、認証、レート制限、障害、応答遅延は Lean の小さなモデルだけでは保証できません。 \\
DB の実挙動 & インデックス、NULL、ロック、トランザクション分離レベル、実行計画は実 DB で確認する必要があります。 \\
UI とブラウザ挙動 & レイアウト、アクセシビリティ、イベント、端末差は実行環境との結びつきが強い領域です。 \\
性能とスケーラビリティ & 時間、メモリ、I/O、負荷時の挙動は測定が必要です。 \\
監視、運用、障害対応 & ログ、アラート、リトライ、フェイルオーバーは運用環境を含むテストが必要です。 \\
\bottomrule
\end{longtable}

通常テストでは、具体例を見られます。
Lean の証明では、モデル化した範囲を全称的に見られます。
どちらが上かではなく、見ている対象が違います。

「Lean で証明したからテストは不要」と述べてはいけません。
Lean で証明したのは、Lean に書いたモデルと定理です。
本番実装がそのモデルと同期していること、外部サービスが契約どおり動くこと、性能要件を満たすことは、別の確認が必要です。

先ほどの料金計算でも、テストは役に立ちます。
たとえば、代表的な入力について期待値を確認することは、仕様レビューや回帰テストとして意味があります。

\begin{listing}[htbp]
\begin{minted}{lean4}
#eval chargeImpl { base := 100, discount := 10, shipping := 5 }  -- => 95

example :
    chargeImpl { base := 100, discount := 10, shipping := 5 } = 95 := by
  rfl

example :
    chargeImpl { base := 8, discount := 10, shipping := 5 } = 5 := by
  rfl
\end{minted}
\caption{\texttt{chargeImpl}}
\label{lst:ch35-charge-examples}
\end{listing}

\texttt{\#eval} は式を評価して結果を表示し、\texttt{example} は名前を付けない定理をカーネルに検査させます。
どちらも、外部 I/O を含む本番コードを実行する通常テストそのものではありません。
二つ目の例では、\texttt{Nat} の引き算がゼロで止まるため、\texttt{8 - 10} は \texttt{0} として扱われます。
これは、現実の料金計算では注意が必要な点です。
このような具体例は、モデルの前提に人間が気づくための確認として有用です。

テストは「この入力ではどうなるか」を明らかにします。
証明は「この形でモデル化したすべての入力で何が成り立つか」を明らかにします。
サンプルテストは、証明の代わりにはなりませんが、モデルの読み方を確認する入口になります。

\section{仕様の過剰形式化を避ける}

Lean 導入でよくある失敗は、最初から仕様全体を形式化しようとすることです。
自然言語仕様には、外部依存、運用判断、業務上の例外、将来変更されるルールが含まれます。
これらをすべて Lean に入れると、証明作業は仕様整理ではなく、巨大な再実装になります。

\textbf{過剰形式化}とは、証明する価値に比べて、モデル化、証明、保守のコストが大きすぎる状態です。
特に、変更頻度が高い周辺仕様、外部システム依存、UI 表示、性能要件を、最初からすべて Lean に入れようとすると起きやすくなります。

過剰形式化を避けるには、次の三つを分けます。

\begin{enumerate}
\item \textbf{証明する核}：Lean theorem として残す部分です。
\item \textbf{テストする周辺}：実装・外部環境・性能・UIなどを確認する部分です。
\item \textbf{文書化する前提}：Lean モデルが何を抽象化し、何を捨てたかを記録する部分です。
\end{enumerate}

\FigurePlaceholder{fig-ch35-proof-kernel.png}{0.90}{証明可能な核}{fig:ch35-proof-kernel}

図\ref{fig:ch35-proof-kernel}のように、実務仕様全体をそのまま Lean に入れるのではなく、中心にある小さな性質を抜き出します。
周辺にある環境依存の振る舞いは、テストや監視に残します。

たとえば「注文金額を正しく計算する」という仕様全体には、商品価格、割引、税、丸め、通貨、決済 API、キャンセル、在庫、監査ログが含まれるかもしれません。
この全体を一度に証明しようとはしません。
最初は、「割引後小計に送料を足す純粋関数が仕様式と一致する」のような核だけを証明します。

過剰形式化を避けるために、設計文書には次のような境界を書くとよいでしょう。

\begin{itemize}
\item この Lean モデルでは、金額を \texttt{Nat} として扱います。
\item 小数、丸め、通貨単位、税率変更はモデル外とします。
\item 外部決済 API の成功と失敗は統合テストで確認します。
\item この theorem は、純粋な計算核の意味保存だけを保証します。
\end{itemize}

このように境界を書くことで、「証明済み」という言葉が過大に解釈されることを防げます。

\section{証明可能な核を抽出する}

証明可能な核は、自然言語仕様から機械的に出てくるわけではありません。
人間が設計判断として切り出す必要があります。
実務で使いやすい手順は、次のとおりです。

\begin{enumerate}
\item 仕様文から、入力、出力、状態、操作を抜き出します。
\item 外部依存を一度外し、純粋な関数として書ける部分を探します。
\item 守りたい性質を、等式、含意、不変条件、round-trip などに言い換えます。
\item モデル化の前提を明示します。
\item theorem 名を仕様項目 ID と対応付けます。
\end{enumerate}

ここでは、出金処理の核を例にします。
現実の出金には、認証、監査ログ、DB 更新、並行性、通知などがあります。
しかし、残高と出金額だけを見る純粋な核は、\texttt{Option Nat} を返す関数として表せます。

\begin{listing}[htbp]
\begin{minted}{lean4}
def withdrawCore (balance amount : Nat) : Option Nat :=
  if _h : amount <= balance then
    some (balance - amount)
  else
    none

theorem withdrawCore_success
    (balance amount : Nat) (h : amount <= balance) :
    withdrawCore balance amount = some (balance - amount) := by
  simp [withdrawCore, h]

theorem withdrawCore_failure
    (balance amount : Nat) (h : ¬ (amount <= balance)) :
    withdrawCore balance amount = none := by
  simp [withdrawCore, h]
\end{minted}
\caption{\texttt{withdrawCore\_success}・\texttt{withdrawCore\_failure}}
\label{lst:ch35-withdraw-core}
\end{listing}

このコードでは、出金可能な場合と不可能な場合を定理として分けています。
\texttt{withdrawCore\_success} は、出金額が残高以下なら、結果は残高から出金額を引いた値になることを述べます。
\texttt{withdrawCore\_failure} は、出金額が残高を超えるなら失敗することを述べます。

ここで重要なのは、DB もログも登場しないことです。
これは欠落ではなく、意図的な抽象化です。
Lean では、まず純粋な残高更新ルールだけを証明します。
DB 更新や監査ログは、別途統合テストや運用テストで確認します。

\subsection{圏論的な読み方}

圏論の言葉に戻すと、証明可能な核を抽出する作業は、複雑な実務世界から、小さな型と関数の世界へ写す作業です。
対象は \texttt{Nat} や \texttt{Option Nat} のような型になり、射は \texttt{withdrawCore} や \texttt{chargeImpl} のような関数になります。

このとき、守りたい性質は、射の等式や保存性として書かれます。

\begin{itemize}
\item 実装が仕様と一致する：\texttt{impl x = spec x}
\item 変換して戻すと元に戻る：\texttt{rollback (migrate x) = x}
\item adapter と map が可換である：\texttt{adapter (map f x) = map f (adapter x)}
\item 操作後も不変条件を満たす：\texttt{Invariant x -> Invariant (step x)}
\end{itemize}

この形まで落とせるなら、Lean で扱う候補になります。
逆に、この形に落とすために大量の外部環境を再現しなければならないなら、通常テストや監視に残す判断が自然です。

\section{設計文書に Lean theorem を埋め込む}

Lean theorem は、Lean ファイルの中だけに閉じ込めると、チームの設計判断に接続しにくくなります。
設計文書には、自然言語の仕様項目と theorem の対応を書いておくとよいでしょう。

推奨する書き方は、「仕様項目 ID」「Lean theorem」「前提条件」「保証すること」「保証しないこと」「対応するテスト」を並べる形式です。
これにより、Lean の証明がどこまでの意味を持つのかをレビューで確認しやすくなります。

たとえば、料金計算の核に対して、次のような設計文書を置けます。

\begin{listing}[htbp]
\begin{minted}{text}
仕様項目: SPEC-CHARGE-001
主張: 実装モデル chargeImpl は仕様モデル chargeSpec と一致する。
Lean theorem: Chapter35.SPEC_CHARGE_001
前提: 金額は Nat。小数、丸め、税率、外部決済APIはモデル外。
保証: 任意の PriceInput について chargeImpl i = chargeSpec i。
保証しないこと: 実DB、外部API、画面表示、性能、通貨丸め。
対応テスト: 決済API統合テスト、丸め仕様テスト、UI表示テスト。
\end{minted}
\caption{\texttt{SPEC-CHARGE-001} 対応表}
\label{lst:ch35-design-doc-fragment}
\end{listing}

この文書断片に対応する theorem を、Lean 側では次のように書けます。

\begin{listing}[htbp]
\begin{minted}{lean4}
/--
SPEC-CHARGE-001:
実装モデル chargeImpl は仕様モデル chargeSpec と一致する。
前提: 金額は Nat。丸め、税率、外部決済APIはモデル外。
-/
theorem SPEC_CHARGE_001 (i : PriceInput) :
    chargeImpl i = chargeSpec i := by
  exact charge_impl_matches_spec i

end Chapter35
\end{minted}
\caption{\texttt{SPEC\_CHARGE\_001}}
\label{lst:ch35-spec-charge-001}
\end{listing}

この対応があると、レビューで次の会話ができます。

\begin{itemize}
\item 仕様項目 \texttt{SPEC-CHARGE-001} は、Lean で証明されているでしょうか。
\item theorem は、どのモデルの上で成り立つでしょうか。
\item モデル外の性質は、どのテストで確認するのでしょうか。
\item 実装変更時にこの theorem が壊れたら、何を見直すのでしょうか。
\end{itemize}

\subsection{定理名を仕様項目に寄せる}

Lean では、数学的に短い theorem 名を付けたくなることがあります。
しかし実務で使うなら、仕様項目 ID との対応が追いやすい名前も有効です。
\texttt{SPEC\_CHARGE\_001} のような名前は数学的には少し無骨ですが、設計文書、Issue、Pull Request、CI ログから検索しやすくなります。

一方、すべての定理を仕様 ID だけにすると、Lean コードが読みにくくなります。
そこで、次のように分けるとよいでしょう。

\begin{itemize}
\item 開発中の補題には、\texttt{charge\_impl\_matches\_spec} のような意味のある名前を付けます。
\item 設計文書に公開する定理には、\texttt{SPEC\_CHARGE\_001} のような仕様 ID 名を付けます。
\item 公開定理は補題を使って短く証明します。
\end{itemize}

この分け方をすると、Lean ファイルが読みやすくなり、設計文書との対応も追いやすくなります。

\section{実装コードと Lean モデルの距離を明示する}

Lean で最も誤解されやすい点は、本番コードと Lean モデルの距離です。
Lean に書いた関数が本番コードそのものではない場合、証明は「本番コードが正しい」ことを直接意味しません。
証明が意味するのは、「Lean に書いたモデルでは、この性質が成り立つ」ということです。

そのため、設計文書には次の情報を書きます。

\begin{itemize}
\item 本番コードのどの関数またはモジュールに対応する Lean モデルでしょうか。
\item 本番コードから何を抽象化したのでしょうか。
\item 本番コードと Lean モデルを同期する責任者またはレビュー観点は何でしょうか。
\item Lean で証明した性質を補うテストは何でしょうか。
\end{itemize}

Lean モデルと本番コードの差分を隠してはいけません。
差分を隠すと、「証明済み」という言葉だけが独り歩きします。
差分を明示すれば、証明は過剰な保証ではなく、設計判断を支える正確な証拠になります。

\section{本章のまとめ}

Lean で証明する価値が高いのは、小さく、純粋で、すべての入力について保証したい性質です。

通常のテストは、外部環境、UI、性能、統合挙動、運用上の振る舞いを確認するために必要です。

過剰形式化を避けるには、証明する核、テストする周辺、文書化する前提を分けます。

自然言語仕様から、型、関数、述語、theorem に落とせる部分を抽出すると、Lean で扱いやすくなります。

設計文書には、Lean theorem、前提条件、保証すること、保証しないこと、対応テストを対応付けて残すとよいでしょう。
